\section{Introduction}
\label{sec:introduction}

Mobile banking is central to retail banks in Portugal. Because these applications run on
untrusted consumer devices and interact directly with the core banking networks, any flaw or
bug in the implementation or design can lead to leaving users exposed to many kinds of attacks.
As such, building a robust banking mobile application is essencial to maintain the security of
data and users. However, security is still often treated as an afterthought. This paper applies
a Secure Software Development Lifecycle (SSDL) to a Portuguese retail bank's mobile banking application, creating a phase by phase
mapping, measurable security requirements, and comparision against secure design principles.

Past failures rarely needed a sophisticated attack. In 2017, Equifax
exposed the data of about 147 million people through an Apache Struts flaw
patched two months earlier, because it had no inventory of where the
component ran and no working patch process~\cite{equifax2018}. In 2012, Knight Capital
lost about 440 million dollars in 45 minutes after a deployment reached
only seven of eight servers~\cite{knight2013}. Both were lifecycle
failures.

\subsection{Security Context}

Mobile banking applications are subject to strict regulatory requirements in Portugal and
the European Union, including the Payment Services Directive~2 (PSD2)~\cite{psd2},
the Digital Operational Resilience Act (DORA)~\cite{dora}, and the General Data
Protection Regulation (GDPR). PSD2 mandates Strong Customer Authentication (SCA) for
all payment transactions, while DORA requires financial entities to demonstrate
operational resilience and incident response capability. Given the high value of
financial assets and sensitive personal data handled, such applications face a broad
threat landscape including credential stuffing, in which the attackers try usernames and
passwords leaked from other services, phishing, where the users are tricked by fraudsters
into revealing credentials or authentication codes through a variaty of methods, such as
fake websites, SMS or phone calls, man-in-the-middle attacks, where an attacker attempts
to intercept or manipulate communications between two parties, allowing attackers to
steal credentials or execute unauthorised actions to, and, finally, API abuse, where
attackers exploit backend weaknesses to do a variaty of different outcomes, such as
steal personal information, execute admin actions and much more.

With the introduction of the Cyber Resilience Act (CRA), most companies, specifically
those that handle critical infrastructure (such as banks) and personal information,
are required to include security into planning, design, development and maintenance
of their digital products. Companies are now mandated, by law, to maintain a Software
Bill of Materials (SBOM), which boils down to being a list of all the third-party
software components and dependencies that the service is dependent on, as well as their
versions. The SBOM helps companies identify quicker vulnerable third-party libraries.
CRA regulation also mandates that companies need to shift from a reactive mindset to a
proactive, continuous monitoring structure, where the companies need to handle
vulnerabilities, detect, manage, patch and report security flaws
in their products during its expected lifecycle. Additionally, CRA dictates that products
must be delivered to the end user with the most secure configurations already activated,
mandating the Secure by Default principle.

Fraudulent or lost transfers are direct losses for the bank, and each
incident damages its reputation. Since developers often neglect security
for lack of resources and priority~\cite{assal2018soups}, and many
vulnerabilities come from predictable human errors~\cite{anu2020human},
security activities must be planned, assigned and verified.

\subsection{Scope and Approach}

The system considered for in this paper is a mobile application for a Portuguese retail
bank. It provides customers with the essencial banking services, including SEPA and
instant bank transfers, MB~WAY payments, card managements (freezing, setting spending
limits, changing PINs), and access to account statements and transaction history.
The application interacts with three main external interfaces: the bank's core banking
system via an internal REST~API, the MB~WAY interoperability network, and the Portuguese
national payment scheme (Multibanco/SIBS). Two categories of users were considered:
customers, who access the banking services and manage their own accounts, and system
administrators, who execute system managing tasks that are subjected to a role based
access control.

The application (i)~binds authentication to the device using biometrics
and a six-digit \ac{PIN} instead of SMS \acp{OTP}; (ii)~ends sessions after
two minutes of inactivity and uses short-lived tokens; (iii)~signs every
sensitive operation, so the server can verify what the user approved;
(iv)~executes transfers as atomic and idempotent operations; and
(v)~rejects requests under overload instead of degrading the isolation
between users.

This paper applies the SSDL to the mobile
banking application, mapping security activities to each development phase following
the NIST Secure Software Development Framework (SSDF, SP~800-218)~\cite{nist800218}
and drawing on OWASP SAMM~v2~\cite{samm} for maturity assessment. The remainder
of the paper is organised as follows: Section~\ref{sec:system} describes the
system and evaluates it against ten secure design principles; Section~\ref{sec:lifecycle}
describes the security activities per lifecycle phase; Section~\ref{sec:cra} presents the
Cyber Resilience Act quick-scan; Section~\ref{sec:conclusion} concludes.
\acused{PSD2}\acused{DORA}\acused{GDPR}\acused{SCA}\acused{CRA}\acused{SBOM}\acused{SSDL}\acused{SSDF}\acused{SAMM}
